\documentclass[11pt]{article}
\usepackage{amsmath,amssymb,amsthm,booktabs,hyperref}
\newtheorem{theorem}{Theorem}
\newtheorem{lemma}[theorem]{Lemma}
\newtheorem{corollary}[theorem]{Corollary}
\title{Two graph classes with three isolated middle vertices}
\author{human\_compression project}
\date{}
\begin{document}\maketitle

\paragraph{Original result and present scope.}
Cambie, Goedgebeur, Jooken and Van den Eede,
\emph{On the order-diameter ratio of girth-diameter cages},
Discrete Mathematics and Theoretical Computer Science 28:4, SOFSEM 2026
\#1, arXiv:2511.21144v3, 19 August 2026, Appendix~B, Table~6
(physical page~23), records order~32 and exactly $25\,408$ ordinary
isomorphism classes for the quintic girth-four, diameter-eight cage.
The final parameter in $(32,5,4,8)$ is a distance, not a uniform
edge--four-cycle multiplicity. No edge-girth-regular hypothesis is used.

This note classifies one explicitly stated profile. It does not classify
every possible middle profile, identify a currently pending computational
family, establish the full $25\,408$ count, or adopt a new checkpoint.
No Source, graph representation, generator or proof checker was executed.

\paragraph{Ambient hypotheses and definitions.}
Let $G$ be a connected simple triangle-free $5$-regular graph on $32$
vertices, of diameter exactly eight. Fix a diametral geodesic
$p_0,p_1,\ldots,p_8$. Put
\[
 A=N(p_0),\quad B=N(p_4),\quad C=N(p_8),\qquad
 I=A\cup B\cup C,\quad J=V(G)\setminus I,\quad F=G[J].
\]
On $K=J\setminus\{p_0,p_4,p_8\}$ define
\[
 L=\{x\in K:N(x)\cap A\ne\varnothing\},\qquad
 R=\{x\in K:N(x)\cap C\ne\varnothing\},\qquad
 M=K\setminus(L\cup R).
\]
The additional hypothesis of the theorem is \emph{exactly}
\[
 |M|=3,\qquad d_F(h)=0\quad\hbox{for every }h\in M.       \tag{H}
\]
In particular, it is not enough to know that three isolated middle
vertices occur if a fourth nonisolated middle vertex is allowed.
Every definition here concerns the complete graph, not merely absent
edges in a partial search state.

\begin{lemma}\label{setup}
The set $I$ is independent, $|I|=15$, $|J|=17$, and $F$ has five edges.
The three centers $p_0,p_4,p_8$ are isolated in $F$. The sets $L,R,M$
partition $K$, and $|L|,|R|\geq5$. There is no edge between $L$ and $R$.
If $D=G[L]$ and $x\in L$, then $d_D(x)\leq |L|-4$; the analogous
inequality holds on $R$.
\end{lemma}
\begin{proof}
Each of $A,B,C$ is independent by triangle-freeness. Two centers are
at distance at least four, so their neighborhoods are disjoint and have
no edge between them: a common neighbor or such an edge would give a
center-to-center path of length two or three. This proves independence
and size~15. All neighbors of a center belong to $I$, so the centers
are isolated in $F$. The graph has $80$ edges, and the independent set
$I$ has $75$ incident edges, all to $J$. Hence $e(F)=80-75=5$.

A vertex meeting both $A$ and $C$ would give a $p_0$--$p_8$ path of
length four. An edge from $L$ to $R$ would similarly give one of length
five. Neither is possible. Each $a\in A$ has, besides $p_0$, four
distinct neighbors in $L$. Thus $|L|\geq4$. Equality would make
$\{p_0\}\cup L$ and $A$ an entire $K_{5,5}$ component: every vertex
of $L$ would already have its five neighbors in $A$. This contradicts
connectedness and order~32. Therefore $|L|\geq5$, and similarly for $R$.

For $x\in L$, choose $a\in A\cap N(x)$. The four $L$-neighbors of $a$
include $x$ and three vertices none of which can neighbor $x$, by
triangle-freeness. Hence $d_D(x)\leq (|L|-1)-3=|L|-4$.
\end{proof}

\begin{lemma}\label{sides}
Under \textup{(H)}, after naming the smaller side $L$,
\[
 |L|=5,\quad |R|=6,\quad |B_L|=3,\quad |B_R|=2,
 \qquad G[L]=K_2\sqcup3K_1,\quad G[R]=C_4\sqcup2K_1.
\]
Here $B_L,B_R$ are the vertices of $B$ whose remaining neighbor lies
in the indicated side.
\end{lemma}
\begin{proof}
An isolated $h\in M$ has all five neighbors in $I$, and meets neither
$A$ nor $C$. Thus $N(h)=B=N(p_4)$. Each $b\in B$ already meets
$p_4$ and the three vertices of $M$, so has precisely one remaining
neighbor, in $L$ or $R$. Both kinds occur: otherwise the opposite
side, its endpoint neighborhood and endpoint center would form a
proper connected component. Thus $B=B_L\sqcup B_R$ with both parts
nonempty.

Since $|K|=14$ and $|M|=3$, the sides have sizes five and six.
Write $b=|B_L|$ and $e=e(G[L])$. The twenty $A$--$L$ incidences,
the $b$ remaining $B_L$--$L$ incidences and the internal edges give
\[
 25=20+b+2e,\qquad 1\leq b\leq4.                         \tag{1}
\]
Lemma~\ref{setup} makes $G[L]$ a matching. The possibilities in (1)
are $(b,e)=(3,1)$ or $(1,2)$. Two matching edges on five vertices
have independence number three, whereas each vertex of $A$ requires
four independent neighbors in $L$. Hence $(b,e)=(3,1)$.
The other side has two $B$-ports and four internal edges, since $F$
has five edges and none touches $M$ or a center or joins the sides.

The graph $D=G[R]$ is triangle-free, has six vertices and four edges,
and has maximum degree two by Lemma~\ref{setup}. Its components
are paths or cycles. If it has a cycle, it is $C_4\sqcup2K_1$.
Otherwise it has two path components, with sizes $(1,5),(2,4)$
or $(3,3)$.

The graph $P_2\sqcup P_4$ has independence number three and cannot
contain the four neighbors required by a vertex of $C$.
In $P_1\sqcup P_5$ the independent set of size four is unique;
in $P_3\sqcup P_3$ it is also unique. It consists of path endpoints
and, in the first case, the isolated vertex and the middle vertex
of $P_5$. Consequently every $C$-vertex would use that same four-set.
At least one nonisolated vertex of $D$ would then have five $C$
neighbors in addition to a positive $D$-degree, contradicting regularity.
Only $C_4\sqcup2K_1$ remains.
\end{proof}

\paragraph{Weighted expansion convention.}
A weighted simple graph assigns a positive integer $w(v)$ to each
vertex. Its expansion replaces $v$ by $w(v)$ independent vertices,
and replaces each edge by the complete bipartite graph between the
corresponding two classes. A zero-weight row in the next table is
omitted. This is a false-twin expansion, not a locally bijective cover.

For $t\in\{1,2,3,4\}$ define $Q_t$ by Table~\ref{quotient}; its edges
are exactly the listed symmetric neighborhood relations. Let $H_t$
be its expansion. Capital and lower-case labels in this table are
vertex names, not assertions about automorphism orbits.

\begin{table}[h]
\centering
\begin{tabular}{ccl}
\toprule
Vertex&Weight&Neighbors\\\midrule
$X$&4&$a_u,a_v$\\
$H$&4&$b_u,b_v,b_p,b_q$\\
$Z$&3&$c_{pq},c_T$\\
$u$&1&$v,a_u,b_u$\\
$v$&1&$u,a_v,b_v$\\
$p$&1&$T,b_p,c_{pq}$\\
$q$&1&$T,b_q,c_{pq}$\\
$T$&2&$p,q,c_T$\\\midrule
$a_u$&$t$&$X,u$\\
$a_v$&$5-t$&$X,v$\\
$b_u$&$4-t$&$H,u$\\
$b_v$&$t-1$&$H,v$\\
$b_p$&1&$H,p$\\
$b_q$&1&$H,q$\\
$c_{pq}$&2&$Z,p,q$\\
$c_T$&3&$Z,T$\\\bottomrule
\end{tabular}
\caption{The complete weighted templates. Zero-weight $b_u$ or $b_v$
and its incident relations are omitted.}\label{quotient}
\end{table}

\begin{theorem}\label{classification}
Exactly two ordinary, unrooted graph isomorphism classes satisfy the
ambient hypotheses and admit a geodesic with profile \textup{(H)}.
Representatives are $H_1$ and $H_2$.
Each representative has order~32, degree~5, girth exactly four and
diameter exactly eight, and is nonbipartite.
\end{theorem}
\begin{proof}[Exhaustion]
Let $uv$ be the sole edge on $L$. Its other three vertices are isolated
in $F$. Every $A$-vertex must choose all three isolates and exactly one
of $u,v$, since it requires four independent neighbors on this five-side.
Let $t$ choose $u$ and $5-t$ choose $v$. Because $u,v\in L$, both
numbers are positive, so $1\leq t\leq4$.

Every one of those three isolates has its five neighbors in $A$,
and is a false twin of $p_0$. They and $p_0$ form $X$ of weight four.
The degree equations at $u,v$ force $4-t$ of the three $B_L$-vertices
to choose $u$, and $t-1$ to choose $v$. This supplies precisely the
classes $a_u,a_v,b_u,b_v$ and their rows in the table.
The center $p_4$ and the three vertices of $M$ form the class $H$ of
weight four, all with neighborhood $B$.

On $R$, each $C$-vertex chooses both isolates and an opposite pair
of the four-cycle. If $k$ choose one opposite pair, $5-k$ choose
the other. Each cycle vertex has two internal neighbors, hence at
most three $C$-neighbors. Thus $2\leq k\leq3$. Exchange the opposite
pairs if necessary to take $k=2$. Denote the pair with two $C$
neighbors by $p,q$ and the other pair by $y_1,y_2$.
Both $p,q$ need one further neighbor in $B_R$; the two $B_R$ vertices
each have just one such neighbor, so these attachments are separate.
The vertices $y_1,y_2$ have neighborhood $\{p,q\}$ together with
the three $C$-vertices choosing their pair. They form $T$ of weight two.
The two isolates on $R$ have neighborhood $C$ and join $p_8$ in the
class $Z$ of weight three. This gives $b_p,b_q,c_{pq},c_T$ and all
remaining relations in the table. Every original vertex and every
possible edge has now been assigned, so $G\cong H_t$.
\end{proof}

\begin{proof}[Existence and exact graph parameters]
In the table the first eight weights sum to $17$, the remaining
weights sum to $15$, and every weighted neighbor sum is five.
Thus $H_t$ is simple and $5$-regular on $32$ vertices. The last
eight rows form an independent set. On the first eight rows the only
edges are $uv$ and the expansion of $pT,qT$, the latter being a
four-cycle. Each last-row neighborhood is independent in this graph.
There is therefore no triangle. Two distinct vertices of $H$ and the
vertices $b_p,b_q$ give a four-cycle. Thus the girth is exactly four.

Both $a_u,a_v$ have positive weight. Picking one vertex from each of
$u,v,a_v,X,a_u$ gives the five-cycle
\[
 u,v,a_v,X,a_u,u,
\]
so the graph is not bipartite.

The graph is connected. On the left, $u,v$ are joined and at least
one of $b_u,b_v$ is present, with total weight three; both $a$ classes
are present. On the right all displayed classes are present and
connected to $H$ through $b_p,b_q$.

Remove the four vertices of $H$. The left classes
$X,u,v,a_u,a_v,b_u,b_v$ and the right classes
$Z,p,q,T,b_p,b_q,c_{pq},c_T$ have no edge between them.
Every vertex on each side has distance at most four to every vertex
of $H$. On the left, $u,v$ have distance at most three, $a_u,a_v$
at most four and $X$ exactly four. On the right the distances are
$1$ for $b_p,b_q$, $2$ for $p,q$, $3$ for $T,c_{pq}$ and $4$ for
$Z,c_T$. In particular $X$ and $Z$ are each at distance exactly four
from $H$, so their mutual distance is eight. All other distances
are at most eight, including those between distinct twins. Hence
$\operatorname{diam}(H_t)=8$.

Choose $p_0$ in $X$, $p_4$ in $H$ and $p_8$ in $Z$. A shortest
left-to-right path through a nonempty $b_u$ or $b_v$ class has length
eight and these centers at positions zero, four and eight. Its three
neighborhoods are exactly the displayed $a$, $b$ and $c$ classes.
The remaining vertices meeting neither endpoint neighborhood are
exactly the other three members of $H$, all isolated in the induced
graph on $J$. Thus the chosen geodesic has profile \textup{(H)}.
\end{proof}

\begin{proof}[Ordinary isomorphism classes]
Exchanging $u,v$, $a_u,a_v$ and $b_u,b_v$ gives a weight-preserving
isomorphism $Q_t\cong Q_{5-t}$. It remains to distinguish $H_1,H_2$
without retaining any root or side colors.

For a simple graph the relation $x\sim y$ defined by $N(x)=N(y)$
is an isomorphism-invariant equivalence relation; distinct equivalent
vertices cannot be adjacent. All vertices in each nonempty table row
are equivalent. For $t=1$ or $t=2$, the listed nonempty rows have
pairwise distinct full neighborhood supports. As the rows represent
disjoint nonempty vertex sets, their expanded neighborhoods are also
pairwise distinct. This comparison includes rows from both halves of
the table. Thus they are precisely the full false-twin classes of the
ordinary expanded graph.

For $t=1$ only $b_v$ is omitted, giving $15$ false-twin classes.
For $t=2$ no row is omitted, giving $16$. These different ordinary
graph invariants prove $H_1\not\cong H_2$. No preservation of $I,J$,
the chosen geodesic, the endpoints or the middle center has been
assumed in this argument. This proves the ordinary two-class claim.
\end{proof}

\begin{corollary}\label{four}
There cannot be four isolated vertices in $M$.
\end{corollary}
\begin{proof}
Each would have neighborhood $B$. Those four vertices and $p_4$
would form one side of a $K_{5,5}$ with $B$, exhausting all degrees
in that component. A connected graph of order~32 cannot contain it.
\end{proof}

\paragraph{A valid whole-family application.}
If \emph{every} target completion of a rooted partial family satisfies
\textup{(H)}, every such completion is one of the two ordinary graphs
above. One can then replace a completion search by finite template
compatibility and explicit transport checks. A transport must respect
every original required and forbidden edge, degree and root/geodesic
constraint and any still applicable history restriction. An unrooted
template match alone does not settle an arbitrary colored partial family.

For example, three fixed vertices already having all five required
neighbors in the fixed set $B$ are permanently isolated in $F$ by
regularity. If every other noncenter vertex of $J$ is forced to meet
$A$ or $C$, then \textup{(H)} follows for every completion. Mere
absence of currently unchosen edges is not such a certificate.

The same ordinary graph may admit several geodesics or middle profiles.
Thus these two classes cannot simply be added to counts from other
strata. A complete census still needs a complete profile cover and
ordinary isomorphism identification across overlaps, or a proved
canonical ownership rule. No division by the number of roots, by a
twin factorial or by an automorphism-group order supplies that rule.

\paragraph{Attribution and acceptance boundary.}
The independent-neighborhood/defect viewpoint, side-incidence methods
and previous isolated-singleton-middle argument already occur in the
retained F009 work. Lemmas~\ref{setup} and~\ref{sides} rederive the
needed parts explicitly. Degree counting, path independence numbers,
false-twin quotients and isomorphism invariants are standard elementary
methods. The present additional assembly is the complete conditional
three-isolated-middle template classification, with no novelty claim.
No original manuscript, pending family, global count, conjecture, Lean
or compression credit is asserted. Actual family matching and distinct
independent mathematical review remain necessary before any adoption.
\end{document}
